For \(g\in\mathcal P_8\), its symbol \(h_g=\{g,\mathcal C_Tg\}/(4\pi)\) is supported in the radius-64 box. At \(K=13\), \(N=47321\), the two sides of the sufficient inequality (ND.24) are the exact integers \[
 N^2=2239277041>935895040=64^2(4+4r_1+2r_2).
\] Therefore (ND.25) and Theorem 1 give a strictly positive quantum averaged-symbol variance for \textbf{every nonzero} such \(g\) at that module, not just the sampled coefficients. The same conclusion holds at every odd K satisfying (ND.24); for each fixed M the condition holds eventually because its right side grows like \(\lambda^{K/2}\) and \(N^2\) grows like \(\lambda^K\).

The actual finite quantum commutator differs from that symbol by a controlled remainder. Put \(G_N=\operatorname{Op}_N(g)\), \(D_N=U_NG_NU_N^*\) and \(R_N=(iN/2)[G_N,D_N]\). The sine remainder gives \[
 \|R_N-\operatorname{Op}_N(h_g)\|_{\rm op}
     \leq\frac{\pi^3}{6N^2}M_3(g)M_3(\mathcal C_Tg),
 \qquad M_3(a)=\sum_v|\widehat a(v)|\,|v|_2^3.
\] Consequently, with normalized Hilbert--Schmidt norm, \[
 \left|\left\|\frac1{\sqrt K}\sum_{j=0}^{K-1}U_N^jR_NU_N^{-j}\right\|_{2,N}
                      -\sqrt{\Gamma_T(g)}\right|
 \leq\frac{\pi^3\sqrt K}{6N^2}M_3(g)M_3(\mathcal C_Tg).      \tag{ND.27}
\] This proves nonvanishing of the averaged commutator for each fixed nonzero radius-eight symbol for all sufficiently large admissible N. It does not assert exact equality between the commutator variance and the classical variance at finite N.

Twelve exact Gaussian-integer coefficient checks compare (ND.25) with a direct modular Weyl-coefficient squared norm, including dimensions 47321 and 1607521. No dense matrices of those sizes are constructed. A small-module negative control, \(K=3,N=7\), violates the non-aliasing condition and gives different finite quantum and classical values; this mismatch is retained rather than removing hypothesis (ND.24).
